\documentclass{amsart}
% For dealing with references we use the comment environment
\usepackage{verbatim}
\newenvironment{reference}{\comment}{\endcomment}
%\newenvironment{reference}{}{}
\newenvironment{slogan}{\comment}{\endcomment}
\newenvironment{history}{\comment}{\endcomment}

% For commutative diagrams we use Xy-pic
\usepackage[all]{xy}

% We use 2cell for 2-commutative diagrams.
\xyoption{2cell}
\UseAllTwocells

% We use multicol for the list of chapters between chapters
\usepackage{multicol}

% This is generally recommended for better output
\usepackage{lmodern}
\usepackage[T1]{fontenc}

% For cross-file-references

% Package for hypertext links:
\usepackage{hyperref}

% For any local file, say "hello.tex" you want to link to please

% Theorem environments.
%
\theoremstyle{plain}
\newtheorem{theorem}[subsection]{Theorem}
\newtheorem{proposition}[subsection]{Proposition}
\newtheorem{lemma}[subsection]{Lemma}

\theoremstyle{definition}
\newtheorem{definition}[subsection]{Definition}
\newtheorem{example}[subsection]{Example}
\newtheorem{exercise}[subsection]{Exercise}
\newtheorem{situation}[subsection]{Situation}

\theoremstyle{remark}
\newtheorem{remark}[subsection]{Remark}
\newtheorem{remarks}[subsection]{Remarks}

\numberwithin{equation}{subsection}

% Macros
%
\def\lim{\mathop{\mathrm{lim}}\nolimits}
\def\colim{\mathop{\mathrm{colim}}\nolimits}
\def\Spec{\mathop{\mathrm{Spec}}}
\def\Hom{\mathop{\mathrm{Hom}}\nolimits}
\def\Ext{\mathop{\mathrm{Ext}}\nolimits}
\def\SheafHom{\mathop{\mathcal{H}\!\mathit{om}}\nolimits}
\def\SheafExt{\mathop{\mathcal{E}\!\mathit{xt}}\nolimits}
\def\Sch{\mathit{Sch}}
\def\Mor{\mathop{\mathrm{Mor}}\nolimits}
\def\Ob{\mathop{\mathrm{Ob}}\nolimits}
\def\Sh{\mathop{\mathit{Sh}}\nolimits}
\def\NL{\mathop{N\!L}\nolimits}
\def\CH{\mathop{\mathrm{CH}}\nolimits}
\def\proetale{{pro\text{-}\acute{e}tale}}
\def\etale{{\acute{e}tale}}
\def\QCoh{\mathit{QCoh}}
\def\Ker{\mathop{\mathrm{Ker}}}
\def\Im{\mathop{\mathrm{Im}}}
\def\Coker{\mathop{\mathrm{Coker}}}
\def\Coim{\mathop{\mathrm{Coim}}}

% Boxtimes
%
\DeclareMathSymbol{\boxtimes}{\mathbin}{AMSa}{"02}

%
% Macros for moduli stacks/spaces
%
\def\QCohstack{\mathcal{QC}\!\mathit{oh}}
\def\Cohstack{\mathcal{C}\!\mathit{oh}}
\def\Spacesstack{\mathcal{S}\!\mathit{paces}}
\def\Quotfunctor{\mathrm{Quot}}
\def\Hilbfunctor{\mathrm{Hilb}}
\def\Curvesstack{\mathcal{C}\!\mathit{urves}}
\def\Polarizedstack{\mathcal{P}\!\mathit{olarized}}
\def\Complexesstack{\mathcal{C}\!\mathit{omplexes}}
% \Pic is the operator that assigns to X its picard group, usage \Pic(X)
% \Picardstack_{X/B} denotes the Picard stack of X over B
% \Picardfunctor_{X/B} denotes the Picard functor of X over B
\def\Pic{\mathop{\mathrm{Pic}}\nolimits}
\def\Picardstack{\mathcal{P}\!\mathit{ic}}
\def\Picardfunctor{\mathrm{Pic}}
\def\Deformationcategory{\mathcal{D}\!\mathit{ef}}

\setlength{\emergencystretch}{3em}
\begin{document}
\title[Stacks Project, Tag 0B1V]{462 --- Intersection of properly meeting cycles on nonsingular projective varieties respects rational equivalence}
\maketitle
\noindent Copyright (C) 2005--2025 Johan de Jong. Stacks Project authors. Source: \href{https://stacks.math.columbia.edu/tag/0B1V}{Tag 0B1V}.

\noindent Editorial difficulty note (not author text). Difficulty H2: The cycles in a rational-equivalence presentation need not meet the fixed cycle properly even when the total cycle does. The proof inducts on excess intersection dimension, uses a moving construction to lower every bad intersection simultaneously, and compares the fibres of a projective one-parameter family while preserving properness at both endpoints..

\noindent Editorial provenance note (not author text). History: excerpt from revision \texttt{a0444\allowbreak 6e57e\allowbreak c1fbc\allowbreak 252a8\allowbreak 71afc\allowbreak ec775\allowbreak 2fb28\allowbreak 07b14}, intersection.tex, lines 2915--3015. Reference presentation was adapted; original-author context and core support blocks were retained or added. The mathematical wording is unchanged. Licensed under GNU FDL 1.2 or later; no invariant sections or cover texts. See ../licenses/COPYING. Context blocks reproduce author definitions and supporting results; they are not additional counted problems.

\section{Conventions}
\label{section-conventions}

\noindent
We fix an algebraically closed ground field $\mathbf{C}$ of any
characteristic. All schemes and varieties are over $\mathbf{C}$ and all
morphisms are over $\mathbf{C}$. A variety $X$ is
{\it nonsingular} if $X$ is a regular scheme (see
Properties, Definition \ref{properties-definition-regular}).
In our case this means that the morphism $X \to \Spec(\mathbf{C})$
is smooth (see
Varieties, Lemma \href{https://stacks.math.columbia.edu/tag/038X}{lemma geometrically regular smooth (Tag 038X)}).




\section{Intersection product using Tor formula}
\label{section-intersection-product}

\noindent
Let $X$ be a nonsingular variety. Let
$\alpha = \sum n_i [W_i]$ be an $r$-cycle and
$\beta = \sum_j m_j [V_j]$ be an $s$-cycle on $X$.
Assume that $\alpha$ and $\beta$ intersect properly, see
Definition \ref{definition-proper-intersection}.
In this case we define
$$
\alpha \cdot \beta = \sum\nolimits_{i,j} n_i m_j W_i \cdot V_j.
$$
where $W_i \cdot V_j$ is as defined in Section \ref{section-tor-formula}.
If $\beta = [V]$ where $V$ is a closed subvariety of dimension $s$,
then we sometimes write $\alpha \cdot \beta = \alpha \cdot V$.



\begin{definition}
\label{properties-definition-regular}
Let $X$ be a scheme. We say $X$ is {\it regular}, or {\it nonsingular} if
for every $x \in X$ there exists an affine open neighbourhood
$U \subset X$ of $x$ such that the ring $\mathcal{O}_X(U)$ is
Noetherian and regular.
\end{definition}

\begin{definition}
\label{definition-proper-intersection}
Let $X$ be a nonsingular variety.
\begin{enumerate}
\item Let $W,V \subset X$ be closed subvarieties with
$\dim(W) = s$ and $\dim(V) = r$. We say that $W$ and $V$
{\it intersect properly} if $\dim(V \cap W) \leq r + s - \dim(X)$.
\item Let $\alpha = \sum n_i [W_i]$ be an $s$-cycle,
and $\beta = \sum_j m_j [V_j]$ be an $r$-cycle on $X$. We say
that $\alpha$ and $\beta$ {\it intersect properly} if
$W_i$ and $V_j$ intersect properly for all $i$ and $j$.
\end{enumerate}
\end{definition}

\section{Intersection multiplicities using Tor formula}
\label{section-tor-formula}

\noindent
A basic fact we will use frequently is that given sheaves of
modules $\mathcal{F}$, $\mathcal{G}$ on a ringed space $(X, \mathcal{O}_X)$
and a point $x \in X$ we have
$$
\text{Tor}_p^{\mathcal{O}_X}(\mathcal{F}, \mathcal{G})_x =
\text{Tor}_p^{\mathcal{O}_{X, x}}(\mathcal{F}_x, \mathcal{G}_x)
$$
as $\mathcal{O}_{X, x}$-modules. This can be seen in several ways
from our construction of derived tensor products in
Cohomology, Section \href{https://stacks.math.columbia.edu/tag/06Y7}{section flat (Tag 06Y7)}, for example it follows from
Cohomology, Lemma \href{https://stacks.math.columbia.edu/tag/06YB}{lemma check K flat stalks (Tag 06YB)}.
Moreover, if $X$ is a scheme and $\mathcal{F}$ and $\mathcal{G}$
are quasi-coherent, then the modules
$\text{Tor}_p^{\mathcal{O}_X}(\mathcal{F}, \mathcal{G})$ are
quasi-coherent too, see
Derived Categories of Schemes, Lemma
\href{https://stacks.math.columbia.edu/tag/08DX}{lemma quasi coherence tensor product (Tag 08DX)}.
More important for our purposes is the following result.



\begin{theorem}
\label{theorem-well-defined}
Let $X$ be a nonsingular projective variety. Let $\alpha$, resp.\ $\beta$
be an $r$, resp.\ $s$ cycle on $X$. Assume that $\alpha$ and $\beta$
intersect properly so that $\alpha \cdot \beta$ is defined. Finally,
assume that $\alpha \sim_{rat} 0$. Then $\alpha \cdot \beta \sim_{rat} 0$.
\end{theorem}

\begin{proof}
Pick a closed immersion $X \subset \mathbf{P}^N$.
By linearity it suffices to prove the result when $\beta = [Z]$ for some
$s$-dimensional closed subvariety $Z \subset X$ which intersects $\alpha$
properly. The condition $\alpha \sim_{rat} 0$ means there
are finitely many $(r + 1)$-dimensional closed subvarieties
$W_i \subset X \times \mathbf{P}^1$ such that
$$
\alpha = \sum [W_{i, a_i}]_r - [W_{i, b_i}]_r
$$
for some pairs of points $a_i, b_i$ of $\mathbf{P}^1$.
Let $W_{i, a_i}^t$ and $W_{i, b_i}^t$ be the irreducible components
of $W_{i, a_i}$ and $W_{i, b_i}$.
We will use induction on the maximum $d$ of the integers
$$
\dim(Z \cap W_{i, a_i}^t),\quad \dim(Z \cap W_{i, b_i}^t)
$$
The main problem in the rest of the proof is that although we know that $Z$
intersects $\alpha$ properly, it may not be the case that $Z$ intersects the
``intermediate'' varieties $W_{i, a_i}^t$
and $W_{i, b_i}^t$ properly, i.e., it may happen that $d >  r + s - \dim(X)$.

\medskip\noindent
Base case: $d = r + s - \dim(X)$. In this case all the intersections of
$Z$ with the $W_{i, a_i}^t$ and $W_{i, b_i}^t$ are proper and the
desired result follows from Lemma \href{https://stacks.math.columbia.edu/tag/0B60}{lemma well defined special case (Tag 0B60)},
because it applies to show that
$[Z] \cdot [W_{i, a_i}]_r \sim_{rat} [Z] \cdot [W_{i, b_i}]_r$ for
each $i$.

\medskip\noindent
Induction step: $d >  r + s - \dim(X)$. Apply
Lemma \href{https://stacks.math.columbia.edu/tag/0B0E}{lemma moving (Tag 0B0E)} to $Z \subset X$ and
the family of subvarieties $\{W_{i, a_i}^t, W_{i, b_i}^t\}$. Then we find a
closed subvariety $C \subset \mathbf{P}^N$ intersecting $X$
properly such that
$$
C \cdot X = [Z] + \sum m_j [Z_j]
$$
and such that
$$
\dim(Z_j \cap W_{i, a_i}^t) \leq \dim(Z \cap W_{i, a_i}^t),\quad
\dim(Z_j \cap W_{i, b_i}^t) \leq \dim(Z \cap W_{i, b_i}^t)
$$
with strict inequality if the right hand side is $> r + s - \dim(X)$.
This implies two things: (a) the induction hypothesis applies to
each $Z_j$, and (b) $C \cdot X$ and $\alpha$ intersect properly (because
$\alpha$ is a linear combination of those $[W_{i, a_i}^t]$ and
$[W_{i, a_i}^t]$ which intersect $Z$ properly).
Next, pick $C' \subset \mathbf{P}^N \times \mathbf{P}^1$
as in Lemma \href{https://stacks.math.columbia.edu/tag/0B1T}{lemma move (Tag 0B1T)} with respect to $C$, $X$, and
$W_{i, a_i}^t$, $W_{i, b_i}^t$.
Write $C' \cdot X \times \mathbf{P}^1 = \sum n_k [E_k]$ for
some subvarieties $E_k \subset X \times \mathbf{P}^1$ of
dimension $s + 1$. Note that $n_k > 0$ for all $k$ by
Proposition \href{https://stacks.math.columbia.edu/tag/0B0V}{proposition positivity (Tag 0B0V)}.
By Lemma \href{https://stacks.math.columbia.edu/tag/0B1M}{lemma transfer (Tag 0B1M)} we have
$$
[Z] + \sum m_j [Z_j] = \sum n_k[E_{k, 0}]_s
$$
Since $E_{k, 0} \subset C \cap X$ we see that $[E_{k, 0}]_s$ and $\alpha$
intersect properly. On the other hand, the cycle
$$
\gamma = \sum n_k[E_{k, \infty}]_s
$$
is supported on $C'_\infty \cap X$ and hence
properly intersects each $W_{i, a_i}^t$, $W_{i, b_i}^t$.
Thus by the base case and linearity, we see that
$$
\gamma \cdot \alpha \sim_{rat} 0
$$
As we have seen that $E_{k, 0}$ and $E_{k, \infty}$
intersect $\alpha$ properly Lemma \href{https://stacks.math.columbia.edu/tag/0B60}{lemma well defined special case (Tag 0B60)}
applied to $E_k \subset X \times \mathbf{P}^1$ and $\alpha$ gives
$$
[E_{k, 0}] \cdot \alpha \sim_{rat} [E_{k, \infty}] \cdot \alpha
$$
Putting everything together we have
\begin{align*}
[Z] \cdot \alpha
& =
(\sum n_k[E_{k, 0}]_r - \sum m_j[Z_j]) \cdot \alpha \\
& \sim_{rat}
\sum n_k [E_{k, 0}] \cdot \alpha \quad (\text{by induction hypothesis})\\
& \sim_{rat}
\sum n_k [E_{k, \infty}] \cdot \alpha \quad (\text{by the lemma})\\
& =
\gamma \cdot \alpha \\
& \sim_{rat}
0 \quad (\text{by base case})
\end{align*}
This finishes the proof.
\end{proof}
\end{document}
